\uuid{TwWY}
\exo7id{6208}
\titre{exo7 6208}
\auteur{queffelec}
\organisation{exo7}
\datecreate{2011-10-16}
\isIndication{false}
\isCorrection{false}
\chapitre{Application linéaire bornée}
\sousChapitre{Application linéaire bornée}
\cpge{mpsi-mp2i.espaces-prehilbertiens.a, pcsi.espaces-prehilbertiens.a, ptsi.geometrie-plan.b, mpsi-mp2i.espaces-prehilbertiens.b, pcsi.espaces-prehilbertiens.b, ptsi.geometrie-espace.b, mpsi-mp2i.applications-lineaires.c, bcpst1.geometrie.f, mp-mpi.topologie-espaces-normes.a, pc.espaces-vectoriels-normes.a, psi.espaces-vectoriels-normes.a, mp-mpi.topologie-espaces-normes.f, psi.espaces-prehilbertiens.e, bcpst2.produit-scalaire.a}
\module{Topologie}
\niveau{L3}
\difficulte{}

\contenu{
\begin{enumerate}
    \item \question{Montrer que pour tout $x\in\Rr^n$, $||x||=\sup\{ |\langle x,y \rangle|\ ; \
||y||\leq1\}$.}
    \item \question{Montrer que l'espace des formes linéaires ${\cal L}(\Rr^n,\Rr)$ sur 
$\Rr^n$ (plus généra\-lement ${\cal L}(E,\Rr)$ où $\dim E$ finie) est
isométriquement isomorphe à $\Rr^n$ (ou $E$).}
\end{enumerate}
}
